\documentclass[11pt,a4paper]{article}
\usepackage[T1]{fontenc}
\usepackage[utf8]{inputenc}
\usepackage{amsmath,amsthm,mathtools}
\usepackage{newtxtext,newtxmath}
\usepackage[a4paper,left=27mm,right=27mm,top=26mm,bottom=26mm,headheight=14pt]{geometry}
\usepackage{microtype,booktabs,array,longtable,needspace}
\usepackage[authoryear,round]{natbib}
\usepackage{enumitem,fancyhdr}
\usepackage[hidelinks,unicode,pdfusetitle]{hyperref}
\usepackage[nameinlink,noabbrev]{cleveref}
\hypersetup{pdftitle={Sharp negative type of finite Ptolemaic metrics},pdfauthor={},pdfsubject={Evidence Press publication manuscript, version 1.0; unrefereed analytic proof and evidence supplement},pdfkeywords={Ptolemaic metric, negative type, Schoenberg matrix, metric inversion, Lorentzian signature}}
\urlstyle{same}
\setlength{\emergencystretch}{1.5em}
\setlength{\parindent}{1.05em}
\linespread{1.025}
\allowdisplaybreaks[1]
\pagestyle{fancy}
\fancyhf{}
\fancyhead[L]{\small Sharp negative type of finite Ptolemaic metrics}
\fancyhead[R]{\small Unrefereed v1.0}
\fancyfoot[C]{\small\thepage}
\renewcommand{\headrulewidth}{0.25pt}
\setlist[enumerate]{label=(\roman*),leftmargin=2em,itemsep=2pt,topsep=4pt}
\newtheorem{theorem}{Theorem}[section]
\newtheorem{lemma}[theorem]{Lemma}
\newtheorem{proposition}[theorem]{Proposition}
\newtheorem{corollary}[theorem]{Corollary}
\theoremstyle{definition}
\newtheorem{definition}[theorem]{Definition}
\newtheorem{example}[theorem]{Example}
\theoremstyle{remark}
\newtheorem{remark}[theorem]{Remark}
\numberwithin{equation}{section}
\newcommand{\one}{\mathbf 1}
\newcommand{\R}{\mathbb R}
\newcommand{\diag}{\operatorname{diag}}
\newcommand{\In}{\operatorname{In}}
\newcommand{\CS}{\operatorname{CS}}
\newcommand{\stp}{\mathsf{S}_p}
\newcommand{\dd}{\mathrm d}
\newcommand{\tr}{^{\mathsf T}}
\newcommand{\abs}[1]{\left|#1\right|}
\newcommand{\norm}[1]{\left\|#1\right\|}
\title{\textbf{Sharp negative type of finite Ptolemaic metrics}}
\author{}
\date{Evidence Press\quad\textbar\quad Version 1.0\\[2pt]25 September 2026}
\begin{document}
\maketitle
\vspace{-1.4em}
\begin{abstract}
For $m\ge3$, let $P(m)$ be the supremum of the exponents $p>0$ for which every $m$-point Ptolemaic metric has $p$-negative type. We prove, for every $m\ge6$, that
\[
P(m)=\log_2\!\left(1+\frac{m}{\lfloor(m-1)^2/4\rfloor}\right),
\]
and that negative type is strict below this exponent. Together with the established three-, four- and five-point cases, this determines $P(m)$ for all $m$ and resolves Conjecture 6.4 of Baker, Huh, Kummer and Lorscheid through their identity $q(n)=P(n-1)$. The proof combines a dimension-independent linear packing estimate for powered star metrics with an induction on normalised entries of inverse augmented distance matrices. Metric inversion allows every principal deletion to remain within the Ptolemaic class. A separate induction on cardinality supplies all positive-definiteness hypotheses from strictly smaller metrics. Complete split metrics give the matching upper bounds. The argument is analytic and imposes no bound on distance ratios; the accompanying exact and symbolic computations check identities and witnesses rather than certify the universal theorem.
\end{abstract}
\noindent\small\textbf{Assurance.} Unrefereed. The proof and supplied review have undergone producer-side checks and replay; proof-assistant verification and external specialist peer review are not claimed.\normalsize

\section{Introduction and main results}\label{sec:intro}
A finite metric space $(X,d)$ is \emph{Ptolemaic} if
\begin{equation}\label{eq:ptolemy}
 d(x,y)d(z,w)\le d(x,z)d(y,w)+d(x,w)d(y,z)
 \qquad(x,y,z,w\in X).
\end{equation}
It has \emph{$p$-negative type} if its entrywise powered distance matrix
$E_p=(d(x,y)^p)_{x,y\in X}$ satisfies
\begin{equation}\label{eq:negative-type}
 z\tr E_pz\le0\qquad\text{whenever }\one\tr z=0.
\end{equation}
The negative type is \emph{strict} when the inequality is strict for every nonzero such $z$. A \emph{snowflake} of a metric $d$ is the metric $d^\alpha$, obtained by raising each distance to a power $0<\alpha\le1$. We use $p$ for the exponent throughout; the notation $q(n)$ is reserved for the triangular-hyperfield threshold discussed below.

Set
\begin{equation}\label{eq:threshold}
 D_m=\left\lfloor\frac{(m-1)^2}{4}\right\rfloor,
 \qquad T_m=1+\frac m{D_m},\qquad p_m=\log_2T_m
 \quad(m\ge3).
\end{equation}
For example, $(T_3,T_4,T_5,T_6)=(4,3,9/4,2)$. The first two newly determined values are
\[
 P(7)=\log_2(16/9),\qquad P(8)=\log_2(5/3).
\]
The sequence $T_m$ is strictly decreasing, and $p_m\le1$ exactly when $m\ge6$.

\begin{theorem}[Sharp threshold for $m\ge6$]\label{thm:main}
For every integer $m\ge6$, every $m$-point Ptolemaic metric has $p_m$-negative type, and has strict $p$-negative type for every $0<p<p_m$. For every $p>p_m$, an explicit complete split metric on $m$ points fails to have $p$-negative type. Consequently,
\[
P(m):=\sup\{p>0:\text{every $m$-point Ptolemaic metric has $p$-negative type}\}
=p_m.
\]
\end{theorem}

The proof also establishes strict $p$-negative type for all $m\ge3$ whenever $0<p<\min\{1,p_m\}$. This weaker small-cardinality statement supplies the induction bases without importing a result at an exponent greater than one.

\begin{corollary}[All cardinalities]\label[corollary]{cor:all}
Together with $P(3)=2$, $P(4)=\log_2 3$ and $P(5)=\log_2(9/4)$, \cref{thm:main} gives
\begin{equation}\label{eq:parity}
 P(2k+1)=2\log_2\frac{k+1}{k}\quad(k\ge1),\qquad
 P(2k)=\log_2\frac{k+1}{k-1}\quad(k\ge2).
\end{equation}
In the notation of Baker, Huh, Kummer and Lorscheid,
\begin{equation}\label{eq:q-formula}
 q(n)=P(n-1)=
 \begin{cases}
 2\log_2\dfrac n{n-2},&n\ge4\text{ even},\\[4pt]
 \log_2\dfrac{n+1}{n-3},&n\ge5\text{ odd}.
 \end{cases}
\end{equation}
\end{corollary}

\subsection{Relation to earlier work}
\citet[Conjecture 6.4 and Proposition 7.5]{bhkl} formulated \eqref{eq:q-formula} in the study of Lorentzian polynomials and weak representations over triangular hyperfields. They identified $q(n)$ with $P(n-1)$ and proved the cases $P(3)$ and $P(4)$. \citet[Theorem 1.1]{ercan} proved $P(5)=\log_2(9/4)$ and $P(6)=1$, using partial correlations, a computer-assisted four-point estimate and coefficient transport under metric inversion. Our proof independently recovers $P(6)$; the newly determined values begin at $P(7)$ and include every larger cardinality. The external five-point endpoint enters only in \cref{cor:all}; it is not used in the proof of \cref{thm:main}.

\citet[Version 3, Sections 11--12]{ksl} identified the general anchor constants and their complete-split thresholds, and established an obstruction above exponent one. That preprint leaves the general low-exponent estimate unresolved (Section 15). Accordingly, the constants and split thresholds themselves are not claimed as new here. The contribution is the universal star-packing estimate and inverse-correlation induction below one, with their strictness hypotheses discharged by the cardinality induction.

The matrix language is classical. Schoenberg's theory identifies negative type with Euclidean embeddability of the corresponding snowflake \citep{schoenberg}. Strict negative type and its exponent range were studied by \citet{liweston}; inverse-distance criteria for the supremal exponent appear in \citet{sanchez}. Augmented distance matrices and their relation to circumcentres belong to distance geometry \citep{fiedler}. We reproduce the identities needed here. The contribution of the present argument is the star-packing estimate in \cref{thm:packing} and its use in the inverse-correlation induction of \cref{thm:inverse-bound}. These replace the need for separate higher-cardinality sign-pattern analyses: packing handles the nonnegative inverse-vector case in every dimension, and a recursive inverse-entry bound handles the remaining signs. Inversion, Schur complements and the complete split witnesses retain the attributions above.

\subsection{Structure of the proof}
For $0<p<1$, the exact star with prescribed distances from a root gives an entrywise lower bound for the Schoenberg matrix of any metric with those root distances. Direct monotonicity of the relevant inverse pairing fails. Instead, a packing theorem controls a nonnegative subsolution of the star's linear system. This proves the required inverse bound whenever either of two inverse vectors is nonnegative.

For the other sign patterns, we select a most negative normalised inverse entry. Deleting a third index either makes it more negative, creates a controlled negative triangle, or forces the nonnegative-vector case. The dimension-dependent constants satisfy the recursion required by this deletion argument.

Finally, a coefficient vector on an $m$-point metric has a smaller sign class of size $a\le m/2$. Every pairwise Schur estimate then involves only $a+2\le m-1$ points. The strict-negative-type induction therefore supplies its hypotheses at the same exponent from strictly smaller metrics. We prove the endpoint by pointwise continuity after the strict induction has closed.

\section{Gram matrices and the four-point base}\label{sec:gram}
Vectors are columns, $\one$ denotes an all-ones vector of the indicated length, and matrix inequalities without $\succeq$ are entrywise. For a symmetric matrix $H$, $H\succ0$ and $H\succeq0$ denote positive definiteness and positive semidefiniteness. We record inertia as the triple of positive, negative and zero eigenvalue counts.

Let $E=(e_{ij})_{i,j=0}^n$ be symmetric with zero diagonal. Take $0$ as a root, put $r_i=e_{0i}$, and define
\begin{equation}\label{eq:gram}
 C_{ij}=\frac{r_i+r_j-e_{ij}}2\quad(1\le i,j\le n).
\end{equation}
In particular, $C_{ii}=r_i$.
If $z_0=-\sum_{i=1}^nz_i$ and $z'=(z_1,\ldots,z_n)\tr$, direct expansion gives
\begin{equation}\label{eq:gram-identity}
 z\tr Ez=-2(z')\tr Cz'.
\end{equation}
Thus $E=d^p$ has strict negative type exactly when $C\succ0$. If $C\succeq0$, choose vectors with Gram matrix $C$ and adjoin the origin; their squared distances are the entries of $E$. When $C\succ0$, these $n+1$ points are affinely independent. This is the finite-dimensional form of Schoenberg's correspondence.

\begin{lemma}[Strictness on at most four points]\label[lemma]{lem:four}
Let $d$ be any metric on at most four points and let $0<p<1$. Then $d$ has strict $p$-negative type.
\end{lemma}
\begin{proof}
The powered array $e_{ij}=d(i,j)^p$ has strict triangle inequalities on distinct points: for $s,t>0$, $(s+t)^p<s^p+t^p$. For four points, write its root Gram matrix as
\[
 C=\begin{pmatrix}s_1&x&y\\x&s_2&z\\y&z&s_3\end{pmatrix}.
\]
The strict triangle inequalities imply
\begin{equation}\label{eq:cut-interval}
 \max\{0,x+y-s_1,x+z-s_2,y+z-s_3\}<\min\{x,y,z\}.
\end{equation}
Indeed, root triangles give $0<x<\min(s_1,s_2)$ and the corresponding inequalities for $y,z$. The three non-root triangle inequalities give, respectively, $x+y-s_1<z$, $x+z-s_2<y$ and $y+z-s_3<x$; the other comparisons follow from the root triangle inequalities.

Choose $\theta$ strictly inside the interval in \eqref{eq:cut-interval}. Then $C$ is the sum of $bb\tr$ over
\[
 b\in\{(1,1,1),(1,1,0),(1,0,1),(0,1,1),(1,0,0),(0,1,0),(0,0,1)\}
\]
with respective coefficients
\[
 \theta,\ x-\theta,\ y-\theta,\ z-\theta,\
 s_1-x-y+\theta,\ s_2-x-z+\theta,\ s_3-y-z+\theta.
\]
All seven coefficients are positive. The last three terms already form a positive definite diagonal matrix, so $C\succ0$.

For three points, $C=\bigl(\begin{smallmatrix}r_1&s\\s&r_2\end{smallmatrix}\bigr)$ with $0<s<\min(r_1,r_2)$, which is positive definite. The two-point case is immediate from \eqref{eq:negative-type}.
\end{proof}

\section{A packing inequality for snowflaked stars}\label{sec:packing}
Fix $0<p<1$ for this section and put
\begin{equation}\label{eq:parameters}
 T=2^p,\qquad \delta=T-1,\qquad a_0=T/2,\qquad c=1-T/2.
\end{equation}
For $n\ge1$, set
\begin{equation}\label{eq:packing-constants}
 R_n(T)=\frac{2n}{2n-(n-1)T},\qquad
 L_n(T)=\max\{2,R_n(T)\}.
\end{equation}
The denominators are positive because $1<T<2$. Both $R_n(T)$ and $L_n(T)$ are non-decreasing in $n$.

Given $\ell_1,\ldots,\ell_n>0$, let $r_i=\ell_i^p$ and let $R$ be the Schoenberg matrix of the powered star with these root-to-leaf lengths:
\begin{equation}\label{eq:star}
 R_{ii}=r_i,\qquad
 R_{ij}=\frac{r_i+r_j-(\ell_i+\ell_j)^p}{2}\quad(i\ne j).
\end{equation}
Put $D=\diag(r)$ and $B=RD^{-1}$. Thus $B_{ii}=1$ and
\begin{equation}\label{eq:b-function}
 B_{ij}=b(\ell_i/\ell_j),\qquad
 b(t)=\frac{1+t^p-(1+t)^p}{2}\quad(i\ne j).
\end{equation}
Strict subadditivity shows that $b(t)>0$, and
\[
 b'(t)=\frac p2\bigl(t^{p-1}-(1+t)^{p-1}\bigr)>0.
\]
Also $b(1)=c$.

\begin{proposition}[A positive-kernel representation]\label[proposition]{prop:kernel}
The matrix $R$ is positive definite in every dimension.
\end{proposition}
\begin{proof}
Substitution and integration by parts give
\[
 x^p=\frac p{\Gamma(1-p)}\int_0^\infty(1-e^{-sx})s^{-1-p}\,\dd s.
\]
Consequently the kernel
\[
 g(x,y)=\frac{x^p+y^p-(x+y)^p}{2}
 =\frac p{2\Gamma(1-p)}\int_0^\infty
 (1-e^{-sx})(1-e^{-sy})s^{-1-p}\,\dd s
\]
is positive semidefinite. The integrals converge: at zero the product is $O(s^2)$, and at infinity it is bounded. Since $g(\ell_i,\ell_i)=cr_i$,
\[
 R=a_0D+[g(\ell_i,\ell_j)]\succ0.
\]
\end{proof}

The following scalar estimate is the main analytic input.
\begin{lemma}[A supporting power]\label[lemma]{lem:scalar}
For $t>0$, define $K(t)=(1+t)^p-t^p$ and $\beta=pc/\delta$. Then
\begin{equation}\label{eq:scalar}
 K(t)\le\delta t^{-\beta},
\end{equation}
with equality exactly when $t=1$.
\end{lemma}
\begin{proof}
Let $h(s)=\log K(e^s)$ and $z=e^s/(1+e^s)\in(0,1)$. Differentiation yields
\begin{align}
 h'(s)&=p\frac{z-z^p}{1-z^p},\label{eq:h-prime}\\
 h''(s)&=\frac{pz(1-z)\left[1-(1-p)z^p-pz^{p-1}\right]}{(1-z^p)^2}<0.
 \label{eq:h-second}
\end{align}
For the sign, weighted arithmetic--geometric mean gives
\[
 (1-p)z^p+pz^{p-1}> (z^p)^{1-p}(z^{p-1})^p=1.
\]
The inequality is strict since the weights are positive and $z^p\ne z^{p-1}$. At zero, $h(0)=\log\delta$ and $h'(0)=-pc/\delta=-\beta$. The graph of the strictly concave function $h$ lies below its tangent at zero, proving \eqref{eq:scalar} and its equality statement.
\end{proof}

\Needspace{6\baselineskip}
\begin{lemma}[Positive solutions of the star system]\label[lemma]{lem:fullsupport}
If $Bw=\one$ and $w>0$, then $\one\tr w\le L_n(T)$.
\end{lemma}
\begin{proof}
Write $M=\sum_iw_i$ and define
\[
 F(t)=\sum_jb(t/\ell_j)w_j.
\]
The function $F$ is increasing. Its diagonal contribution at $t=\ell_i$ is $cw_i$, whereas the diagonal contribution in $Bw$ is $w_i$. Thus
\[
 F(\ell_i)+a_0w_i=1.
\]
It follows that $w_i$ is non-increasing as $\ell_i$ increases.

The $i$th row equation also reads
\[
 M+w_i=2+\sum_{j\ne i}K(\ell_i/\ell_j)w_j.
\]
Set $t_i=\ell_i^{-\beta}$, $H=\sum_iw_i/t_i$ and $S=\sum_i1/t_i$. By \cref{lem:scalar},
\[
 M+Tw_i-\delta t_iH\le2.
\]
After division by $t_i$ and summation, this becomes
\begin{equation}\label{eq:row-sum}
 (M-2)S\le\bigl[(n-1)\delta-1\bigr]H.
\end{equation}
If $M\le2$, the conclusion follows. Otherwise the coefficient on the right is positive. The sequences $w_i$ and $1/t_i=\ell_i^\beta$ are oppositely ordered, so
\begin{equation}\label{eq:chebyshev}
 MS-nH=-\sum_{i<j}(w_i-w_j)(\ell_i^\beta-\ell_j^\beta)\ge0.
\end{equation}
Hence $H/S\le M/n$. Substituting in \eqref{eq:row-sum} gives
\[
 \bigl[n+1-(n-1)\delta\bigr]M\le2n,
\]
which is $M\le R_n(T)$.
\end{proof}

\begin{theorem}[Star packing]\label{thm:packing}
For every positive length vector, every $w\ge0$ satisfying $Bw\le\one$ obeys
\begin{equation}\label{eq:packing}
 \one\tr w\le L_n(T)=\max\left\{2,\frac{2n}{2n-(n-1)T}\right\}.
\end{equation}
When $R_n(T)\ge2$, this bound is attained for equal lengths.
\end{theorem}
\begin{proof}
We induct on $n$; the case $n=1$ is immediate. The polytope
\[
 \mathcal P=\{w\in\R^n:w\ge0,\ Bw\le\one\}
\]
is nonempty and compact. Indeed $B$ has positive entries and diagonal one, so every feasible coordinate lies in $[0,1]$. Choose a vertex maximising $\one\tr w$.

If its support has $k<n$ elements, the corresponding $k$ row inequalities give the same packing problem for the length vector on that support. Induction gives $\one\tr w\le L_k(T)\le L_n(T)$; the omitted row constraints can only reduce the feasible set.

Otherwise all coordinates are positive and none of the nonnegativity constraints is active. A vertex in $\R^n$ requires $n$ linearly independent active constraints. There are exactly $n$ remaining constraints, namely the rows of $Bw\le\one$, so all must be equalities. \Cref{lem:fullsupport} applies.
\end{proof}

\begin{remark}\label[remark]{rem:packing-sharp}
When $R_n(T)\ge2$, equal lengths and the vector
$w=[1+(n-1)c]^{-1}\one$ attain \eqref{eq:packing}. We do not assert that the constant $2$ is optimal in the complementary parameter range. In particular, the theorem does not assert nonnegativity of $R^{-1}\one$ or $R^{-1}r$.
\end{remark}

\section{Inverse pairings, augmented matrices and inversion}\label{sec:inverse}
Fix $0<p<1$ and let $E=d^p$ be a powered metric with a positive definite root Gram matrix $C$ of order $n$. Define
\begin{equation}\label{eq:inverse-quantities}
 r=\diag C,\quad u=C^{-1}r,\quad v=C^{-1}\one,\quad
 \mu=\one\tr u=r\tr v,\quad
 U=\one\tr C^{-1}\one,\quad V=r\tr C^{-1}r.
\end{equation}
Both $U,V$ are positive. Cauchy--Schwarz in the $C^{-1}$ inner product gives
\begin{equation}\label{eq:cs}
 UV\ge\mu^2.
\end{equation}
The root-dependent quantity $\mu$ is only an auxiliary mixed inverse pairing. We do not identify it with the usual metric $M$-constant or claim a new invariant.

\begin{lemma}[Nonnegative inverse vectors]\label[lemma]{lem:nonnegative}
If either $u\ge0$ or $v\ge0$, then $\mu\le L_n(T)$.
\end{lemma}
\begin{proof}
Choose $\ell_i=d(0,i)$ and form the star matrix $R$ in \eqref{eq:star}. The triangle inequality gives $C\ge R$ entrywise, with identical diagonals. If $v\ge0$, then $w=Dv\ge0$ and
\[
 Bw=Rv\le Cv=\one.
\]
\Cref{thm:packing} gives $\mu=\sum_iw_i\le L_n(T)$.

If $u\ge0$, put $C^*=D^{-1}CD^{-1}$ and $R^*=D^{-1}RD^{-1}$. Their common diagonal is $r^*=1/r$, componentwise, and $R^*$ is exactly the star matrix for the reciprocal lengths $1/\ell_i$. Also
\[
 (C^*)^{-1}\one=Du\ge0,\qquad
 (r^*)\tr(C^*)^{-1}\one=\mu.
\]
The preceding packing argument applies to $C^*\ge R^*$. This step is a diagonal congruence calculation and needs no claim of monotonicity of inverse matrices or inverse pairings.
\end{proof}

\subsection{The augmented inverse}
Adjoin an index $\infty$ and set
\begin{equation}\label{eq:augmented}
 A_{\infty i}=A_{i\infty}=1,\quad A_{ij}=e_{ij},\quad A_{ii}=0.
\end{equation}
In the order $\infty,0,1,\ldots,n$,
\[
 A=\begin{pmatrix}
 0&1&\one\tr\\
 1&0&r\tr\\
 \one&r&r\one\tr+\one r\tr-2C
 \end{pmatrix}.
\]
The Schur complement of the leading $2\times2$ block is $-2C$. Thus
\begin{equation}\label{eq:aug-inertia}
 \In(A)=(1,n+1,0),
\end{equation}
and direct block inversion gives
\begin{equation}\label{eq:aug-inverse}
 A^{-1}=\begin{pmatrix}
 -V/2&1-\mu/2&u\tr/2\\
 1-\mu/2&-U/2&v\tr/2\\
 u/2&v/2&-C^{-1}/2
 \end{pmatrix}.
\end{equation}
Whenever the inverse diagonal is negative, write
\begin{equation}\label{eq:rho}
 \rho_{ij}(A)=\frac{(A^{-1})_{ij}}
 {\sqrt{[-(A^{-1})_{ii}][-(A^{-1})_{jj}]}}\quad(i\ne j).
\end{equation}
In particular,
\begin{equation}\label{eq:rho-mu}
 \rho_{\infty0}(A)=\frac{2-\mu}{\sqrt{UV}}.
\end{equation}
Define
\begin{equation}\label{eq:kn}
 k_n(T)=1-\frac2{L_n(T)}
 =\max\left\{0,\frac{(n-1)T-n}{n}\right\}.
\end{equation}
By \cref{lem:nonnegative}, either $u\ge0$ or $v\ge0$ implies
\begin{equation}\label{eq:good-sign}
 \rho_{\infty0}(A)\ge-k_n(T).
\end{equation}
Indeed, for $\mu\le2$ the numerator in \eqref{eq:rho-mu} is nonnegative; for $2<\mu\le L_n(T)$, \eqref{eq:cs} gives
\[
 \rho_{\infty0}(A)\ge\frac{2-\mu}{\mu}\ge\frac2{L_n(T)}-1.
\]

\begin{remark}[Classical distance-geometric quantities]\label[remark]{rem:classical}
If vectors with Gram matrix $C$ are assembled as the columns of $X$, the sphere through these vectors and the origin has centre $XC^{-1}r/2$. Its barycentric coefficient at the root is $1-\mu/2$. This explains the mixed pairing in \eqref{eq:inverse-quantities}. Also $\one\tr E^{-1}\one=2/V$, whenever $C\succ0$: the latter condition gives $n$ negative eigenvalues of $E$ by \eqref{eq:gram-identity} and a positive eigenvalue because its off-diagonal entries are positive. Comparison of the $\infty\infty$ block of \eqref{eq:aug-inverse} with the Schur complement of $E$ proves the formula. These relations place our auxiliary quantities within the classical framework of \citet{fiedler} and the inverse-distance criteria discussed by \citet{sanchez}.
\end{remark}

\subsection{Changing the distinguished index}
A positive diagonal congruence $A\mapsto SAS$ sends its inverse to $S^{-1}A^{-1}S^{-1}$. It preserves all signs and all quantities \eqref{eq:rho}. Any row can be normalised to have off-diagonal entries one.

\begin{lemma}[Inversion closure]\label[lemma]{lem:inversion}
Let $(X,d)$ be Ptolemaic and $o\in X$. On $(X\setminus\{o\})\cup\{\infty\}$ define
\begin{equation}\label{eq:metric-inversion}
 \widehat d(i,j)=\frac{d(i,j)}{d(o,i)d(o,j)},\qquad
 \widehat d(\infty,i)=\frac1{d(o,i)}.
\end{equation}
Then $\widehat d$ is a Ptolemaic metric. Normalising row $o$ of the augmented powered-distance matrix gives the augmented matrix of $\widehat d^{\,p}$.
\end{lemma}
\begin{proof}
All distances between distinct points are positive. Multiplying a triangle inequality among three finite points by the product of their three distances to $o$ gives a Ptolemy inequality for $d$. A triangle inequality involving $\infty$ becomes a triangle inequality involving $o$ in the original metric. Ptolemy inequalities on four finite points are preserved after multiplication by their common denominator. A Ptolemy inequality containing $\infty$ reduces to an ordinary triangle inequality for $d$. This checks every case. For the matrix statement, take the diagonal scale at $o$ to be one and the scale at any other index $i$ to be $1/A_{oi}$.
\end{proof}

Every principal submatrix of an augmented array, of order at least three, is therefore diagonally congruent to an augmented powered-distance matrix of a Ptolemaic metric with one fewer point than its matrix order. If the principal submatrix does not contain the original $\infty$, first normalise one of its indices and then restrict the inverted metric. An order-two principal matrix has inertia $(1,1,0)$ directly.

\subsection{Proper-subset inertia}
For $j\ge2$, let $\stp(j)$ denote the following \emph{universal} statement:
\begin{quote}
Every Ptolemaic metric on at most $j$ points has strict $p$-negative type.
\end{quote}
This includes all inverted metrics, not just subsets of one fixed metric.

\begin{lemma}[Inverse two-by-two blocks]\label[lemma]{lem:pair-inertia}
Assume $\stp(n+1)$ for some $n\ge2$, and let $A$ augment an $(n+1)$-point Ptolemaic metric powered by $p$. Every principal $2\times2$ block of $A^{-1}$ is negative definite. In particular all inverse diagonal entries are negative and $|\rho_{ij}(A)|<1$.
\end{lemma}
\begin{proof}
By \eqref{eq:aug-inertia}, $A$ has inertia $(1,n+1,0)$. Delete a pair of indices $J$ and call the complementary principal matrix $H$. It has order $n$. If $n=2$, $H$ has inertia $(1,1,0)$ directly. Otherwise \cref{lem:inversion} and $\stp(n+1)$ give $\In(H)=(1,n-1,0)$.

The standard block congruence underlying inertia additivity \citep{haynsworth} is
\[
 \begin{pmatrix}I&-H^{-1}Z\\0&I\end{pmatrix}\tr
 \begin{pmatrix}H&Z\\Z\tr&Y\end{pmatrix}
 \begin{pmatrix}I&-H^{-1}Z\\0&I\end{pmatrix}
 =\begin{pmatrix}H&0\\0&Y-Z\tr H^{-1}Z\end{pmatrix}.
\]
Its $2\times2$ Schur complement has inertia $(0,2,0)$. Its inverse is $(A^{-1})_{JJ}$, which is negative definite. The stated diagonal and correlation conclusions follow from its principal minors.
\end{proof}

\section{A normalised inverse-correlation bound}\label{sec:correlation}
The next theorem is conditional on strictness of smaller-cardinality metrics. That hypothesis is discharged only in \cref{sec:cardinality}.

\begin{theorem}[Inverse-correlation induction]\label{thm:inverse-bound}
Fix $0<p<1$, put $T=2^p$, and let $n\ge2$. Assume $\stp(n+1)$. If $n\ge3$, suppose also that
\begin{equation}\label{eq:inverse-range}
 T\le\frac n{n-2}.
\end{equation}
Then every augmented powered-distance matrix $A$ of an $(n+1)$-point Ptolemaic metric satisfies
\begin{equation}\label{eq:inverse-bound}
 \rho_{ij}(A)\ge-k_n(T)\qquad(i\ne j).
\end{equation}
For $n=2$, the conclusion holds without \eqref{eq:inverse-range}.
\end{theorem}
\noindent In the proof for $m$ points, this theorem is applied with $n+1=a+2\le m-1$, so its strictness hypothesis comes entirely from the outer induction.
\begin{proof}
We induct on $n$. When $n=2$, normalise either index of the pair as $\infty$ and take the other as the root. The resulting three-point powered metric has strict triangle inequalities. Its root Gram matrix
$C=\bigl(\begin{smallmatrix}r_1&s\\s&r_2\end{smallmatrix}\bigr)$ has $0<s<\min(r_1,r_2)$, which gives $u,v>0$ directly. Since $k_2(T)=0$, \eqref{eq:good-sign} proves the claim in every gauge.

Now let $n\ge3$. The lower-dimensional strictness hypothesis follows from $\stp(n+1)$, and its parameter condition, when needed, follows because $n/(n-2)$ decreases with $n$. Write $W=A^{-1}$. If no off-diagonal entry is negative, there is nothing to prove. Otherwise choose a most negative normalised entry, at indices $i,j$, and write
\[
 \gamma=-\rho_{ij}(A)>0.
\]
By \cref{lem:pair-inertia}, $\gamma<1$. Suppose, for a contradiction, that $\gamma>k_n(T)$.

For any third index $\ell$, set
\[
 s=-\frac{W_{i\ell}}{\sqrt{(-W_{ii})(-W_{\ell\ell})}},\qquad
 t=-\frac{W_{j\ell}}{\sqrt{(-W_{jj})(-W_{\ell\ell})}}.
\]
Deleting $\ell$ from $A$ changes the inverse by
\begin{equation}\label{eq:inverse-deletion}
 W^{(\ell)}=W_{\bar\ell,\bar\ell}
 -\frac{W_{\bar\ell,\ell}W_{\ell,\bar\ell}}{W_{\ell\ell}}.
\end{equation}
The signed negative magnitude of the new $ij$ entry is therefore
\begin{equation}\label{eq:partial-correlation}
 \gamma'=\frac{\gamma-st}{\sqrt{(1-s^2)(1-t^2)}}.
\end{equation}
The denominator is positive by \cref{lem:pair-inertia}. The deleted matrix can be normalised as the augmentation of an $n$-point Ptolemaic metric, even when $\ell$ is the original $\infty$. The induction hypothesis applies to it.

If $st\le0$, then $\gamma'\ge\gamma>k_n(T)\ge k_{n-1}(T)$, a contradiction. Thus, for each third index, both incident entries are nonzero and have the same sign.

Suppose that for some $\ell$ both incident inverse entries are negative. Then $0<s,t\le\gamma$ by maximality. Since
\begin{equation}\label{eq:denominator}
 (1-st)^2-(1-s^2)(1-t^2)=(s-t)^2\ge0,
\end{equation}
and $\gamma-st\ge\gamma-\gamma^2>0$, we have
\begin{equation}\label{eq:gamma-step}
 \gamma'\ge\frac{\gamma-st}{1-st}
 \ge\frac{\gamma-\gamma^2}{1-\gamma^2}
 =\frac\gamma{1+\gamma}.
\end{equation}
The middle inequality uses $st\le\gamma^2$ and the fact that $(\gamma-x)/(1-x)$ decreases with $x$ for $\gamma<1$.

The constants satisfy
\begin{equation}\label{eq:compatibility}
 k_{n-1}(T)\le\frac{k_n(T)}{1+k_n(T)}.
\end{equation}
If $k_{n-1}=0$ this is immediate; if $k_n=0$, monotonicity in $n$ gives $k_{n-1}=0$ as well. When both are positive, direct simplification yields
\begin{equation}\label{eq:compatibility-factor}
 \frac{k_n(T)}{1+k_n(T)}-k_{n-1}(T)
 =\frac{(T-1)\,[n-(n-2)T]}{T(n-1)}\ge0.
\end{equation}
Since $x/(1+x)$ is increasing, \eqref{eq:gamma-step} and $\gamma>k_n$ contradict the lower-dimensional bound.

The only remaining case is that every third index has two positive incident inverse entries. Normalise $i$ as $\infty$ and take $j$ as the root. In \eqref{eq:aug-inverse} this says $u,v>0$. The metric remains Ptolemaic by \cref{lem:inversion}, and its Gram matrix is positive definite by $\stp(n+1)$. Equation~\eqref{eq:good-sign} bounds the selected entry below by $-k_n(T)$, the final contradiction.
\end{proof}

\subsection{Conditioned correlations}
Let a Ptolemaic metric have $a$ distinguished anchors and two further points. Choose one anchor as root. Assuming positive definiteness, partition its Gram matrix as
\[
 C=\begin{pmatrix}K&L\\L\tr&H\end{pmatrix},\qquad
 S=H-L\tr K^{-1}L,
\]
where $K$ is indexed by the other $a-1$ anchors and $S$ by the two poles. The corresponding pole block of the augmented inverse is $-S^{-1}/2$. For a positive definite $2\times2$ matrix,
\begin{equation}\label{eq:schur-correlation}
 \frac{(-S^{-1})_{12}}
 {\sqrt{(-S^{-1})_{11}(-S^{-1})_{22}}}
 =\frac{S_{12}}{\sqrt{S_{11}S_{22}}}.
\end{equation}
This follows by writing out the two-by-two inverse; in particular the sign has not been lost by squaring.

\begin{corollary}[Anchor--pole estimate]\label[corollary]{cor:anchors}
Suppose $a\ge2$, $\stp(a+2)$ holds, and $T\le(a+1)/(a-1)$. Then
\begin{equation}\label{eq:anchor-bound}
 S_{12}\ge-\kappa_a(T)\sqrt{S_{11}S_{22}},\qquad
 \kappa_a(T)=\max\left\{0,\frac{aT-a-1}{a+1}\right\}.
\end{equation}
\end{corollary}
\begin{proof}
Apply \cref{thm:inverse-bound} with $n=a+1$ and use \eqref{eq:schur-correlation}.
\end{proof}

\section{Induction on the number of points}\label{sec:cardinality}
\begin{proposition}[Strict subcritical negative type]\label[proposition]{prop:strict}
For every $m\ge3$, every $m$-point Ptolemaic metric has strict $p$-negative type whenever
\begin{equation}\label{eq:strict-range}
 0<p<\min\{1,p_m\}.
\end{equation}
Every two-point metric has strict $p$-negative type for every $p>0$.
\end{proposition}
\begin{proof}
We use strong induction on $m$. The cases $m\le4$ follow from \cref{lem:four}. Fix $m\ge5$ and an exponent in \eqref{eq:strict-range}. The sequence $T_j$ decreases with $j$: writing $j=2k+1$ or $2k$ gives
\[
 T_{2k+1}=\left(\frac{k+1}{k}\right)^2,
 \qquad T_{2k}=\frac{k+1}{k-1},
\]
and the two adjacent comparisons are strict. Hence the strong induction hypothesis gives $\stp(m-1)$ at the \emph{current exponent} $p$. In particular it covers all smaller Ptolemaic metrics obtained by inversion and restriction.

Let $z\ne0$ satisfy $\one\tr z=0$. If some coordinates vanish, apply the induction hypothesis to its proper support. Otherwise reverse its sign if necessary, so that it has $a$ positive and $r$ negative coordinates with $a\le r$ and $a+r=m$. By homogeneity, normalise both absolute coefficient sums to one.

If $a=1$, write the negative coefficients as $-\beta_i$, with $\beta_i>0$ and $\sum_i\beta_i=1$. The powered array $E=d^p$ is an ordinary metric, so
\begin{align*}
 -\frac12z\tr Ez
 &=\sum_i\beta_i e_{0i}-\sum_{i<j}\beta_i\beta_j e_{ij}\\
 &\ge\sum_i\beta_i e_{0i}-\sum_{i<j}\beta_i\beta_j(e_{0i}+e_{0j})
 =\sum_i\beta_i^2e_{0i}>0.
\end{align*}

Suppose now that $a\ge2$. Since $a\le r$ and $m\ge5$, we have $r\ge3$ and therefore
\begin{equation}\label{eq:size-guard}
 a+2\le m-1.
\end{equation}
For this sign pattern define
\begin{equation}\label{eq:pattern-threshold}
 F(a,r)=\frac{(a+1)r}{a(r-1)}
 =1+\frac{m}{a(m-1-a)}.
\end{equation}
Since $a(m-1-a)\le D_m$,
\begin{equation}\label{eq:parameter-chain}
 T<T_m\le F(a,r)\le\frac{a+1}{a-1}.
\end{equation}
The last inequality is equivalent to $r\ge a$. Also
\begin{equation}\label{eq:copositive-margin}
 (r-1)\kappa_a(T)<1.
\end{equation}
If the positive part in $\kappa_a$ is inactive this is immediate. Otherwise the exact identity is
\[
 1-(r-1)\kappa_a(T)
 =\frac{a(r-1)}{a+1}\,[F(a,r)-T]>0.
\]

Use one positive point as the Gram root and partition the full Gram matrix as
\[
 C=\begin{pmatrix}K&L\\L\tr&H\end{pmatrix},
 \qquad S=H-L\tr K^{-1}L,
\]
where $K$ is indexed by the remaining $a-1$ positive points. We do \emph{not} assume that $C$ is positive definite. The anchor metric has $a$ points, so $K\succ0$ follows from $\stp(m-1)$. Adding a single negative point uses $a+1$ points and proves $S_{ii}>0$. Adding a pair uses $a+2$ points. By \eqref{eq:size-guard} and \eqref{eq:parameter-chain}, \cref{cor:anchors} applies to each such pair at the same exponent, giving
\[
 S_{ij}\ge-\kappa_a(T)\sqrt{S_{ii}S_{jj}}.
\]
These pairwise Schur complements are precisely the corresponding principal two-by-two blocks of $S$, since they eliminate the same anchor block $K$.

Let $\beta_i>0$ be the negative coefficient magnitudes and put $t_i=\beta_i\sqrt{S_{ii}}$. Then
\begin{align}
 \beta\tr S\beta
 &\ge\sum_i t_i^2-2\kappa_a(T)\sum_{i<j}t_it_j\notag\\
 &=\bigl[1-(r-1)\kappa_a(T)\bigr]\sum_i t_i^2
   +\kappa_a(T)\sum_{i<j}(t_i-t_j)^2
 >0.\label{eq:copositive}
\end{align}
If $\alpha$ is the vector of positive coefficients excluding the root, completion of the square gives
\begin{equation}\label{eq:completion}
 -\frac12z\tr Ez
 =(\alpha-K^{-1}L\beta)\tr K(\alpha-K^{-1}L\beta)
   +\beta\tr S\beta>0.
\end{equation}
This closes the induction. Balanced patterns are included: when $a=r$, the same strict subset bound \eqref{eq:size-guard} holds.
\end{proof}

\begin{proof}[Endpoint part of \cref{thm:main}]
For $m\ge6$, $p_m\le1$. Fix an $m$-point Ptolemaic metric and a zero-sum vector $z$. By \cref{prop:strict}, $z\tr d^pz\le0$ for $0<p<p_m$. This finite sum is continuous in $p$, so letting $p\uparrow p_m$ proves the endpoint inequality. The metric and vector were arbitrary. No inverse is taken at the endpoint, and no compactness assumption on the class of metrics is needed. The matching upper bound follows in the next section.
\end{proof}

\section{Sharpness and the Lorentzian threshold}\label{sec:sharp}
\begin{proposition}[Complete split witnesses]\label[proposition]{prop:sharp}
For every $m\ge3$, there is a Ptolemaic metric on $m$ points that fails $p$-negative type for every $p>p_m$.
\end{proposition}
\begin{proof}
Let $a,r\ge1$, $r\ge2$, and $a+r=m$. The complete split graph consists of a clique of size $a$ and an independent set of size $r$, with every edge between the two parts present. On its vertex set, define the shortest-path metric $\CS(a,r)$ to have all distinct distances one, except distances inside the $r$-point set, which are two. Triangle inequalities hold. For four distinct points, two points in each part give opposite-edge products $2,1,1$; all other distributions give three equal products. Thus Ptolemy holds. Repeated-index inequalities are automatic.

Give every point in the $a$-part coefficient $1/a$ and every point in the $r$-part coefficient $-1/r$. The negative-type inequality becomes
\begin{equation}\label{eq:split-inequality}
 \frac{a-1}{2a}+\frac{r-1}{2r}2^p\le1,
\end{equation}
which is equivalent to $2^p\le F(a,r)$. Choose $a$ to maximise the denominator $a(m-1-a)$ in \eqref{eq:pattern-threshold}. Its integer maximum is $D_m$.

For $m=2k+1$, the choice $(a,r)=(k,k+1)$ gives $T_m=((k+1)/k)^2$. For $m=2k$, the choices $(k-1,k+1)$ and $(k,k)$ both give $T_m=(k+1)/(k-1)$. Hence \eqref{eq:split-inequality} fails for every $p>p_m$.
\end{proof}

\Cref{prop:sharp} and the endpoint argument complete \cref{thm:main}. The three-point equality $P(3)=2$ is also elementary: every three-point metric is a Euclidean triangle, and $\CS(1,2)$ has exponent at most two. The equality $P(4)=\log_2 3$ is \citet[Theorem 7.16 and Corollary 7.17]{bhkl}, while $P(5)=\log_2(9/4)$ is \citet[Theorem 1.1]{ercan}. Appending these two results and the three-point observation proves \eqref{eq:parity}. The entire argument for $m\ge6$ is independent of these imported endpoint results.

\paragraph{Seven-point example.}
For $\CS(3,4)$, give the three clique vertices coefficient $1/3$ and the four independent-set vertices coefficient $-1/4$. With $T=2^p$, direct expansion gives
\[
 z\tr d^p z=\frac23+\frac34T-2.
\]
This expression vanishes at $T=16/9$ and is positive for every larger $T$. It displays the sharpness of the first new threshold, $P(7)=\log_2(16/9)$, without a numerical eigenvalue calculation.

\subsection{The matrix formulation}
We recall the correspondence of \citet[Proposition 7.5]{bhkl} to identify exactly which Lorentzian statement follows. An \emph{admissible} $N\times N$ matrix is symmetric, has zero diagonal and has positive off-diagonal entries. It satisfies the $\mathbb T_p$-inequalities when, for every quartet of distinct indices, the three opposite-edge products raised to the power $1/p$ satisfy the triangle inequalities. Let $q(N)$ be the supremum of the parameters $p$ for which every such matrix has at most one positive eigenvalue.

Normalise one row of an admissible matrix to have off-diagonal entries one, and write it as
\[
 A=\begin{pmatrix}0&\one\tr\\\one&E\end{pmatrix}.
\]
The $\mathbb T_p$-inequalities containing the distinguished index say that $d_{ij}=e_{ij}^{1/p}$ is a metric. The remaining quartet inequalities say that $d$ is Ptolemaic. Conversely, augmenting $d^p$ for a Ptolemaic metric gives exactly these quartet inequalities. The block congruence in \cref{sec:inverse}, which is valid for arbitrary symmetric $C$, shows that $A$ has at most one positive eigenvalue exactly when the root Gram matrix is positive semidefinite. Therefore $q(N)=P(N-1)$. This recovers \eqref{eq:q-formula} and establishes \cref{cor:all}.

\subsection{Embedding and asymptotic consequences}
For $m\ge6$ and $0<p<p_m$, \eqref{eq:gram-identity} gives an affinely independent isometric embedding of $(X,d^{p/2})$ into $\R^{m-1}$. At $p=p_m$, an isometric embedding still exists but may lose affine independence. Strictness below the endpoint for the three smaller cardinalities also follows from strict descent \citep{liweston}; that theorem is not used in our cardinality induction.

The explicit formula gives
\[
 \lim_{m\to\infty}mP(m)=\frac4{\log 2}.
\]
Indeed $D_m=m^2/4+O(m)$, so $T_m=1+4/m+O(m^{-2})$, and $\log(1+x)=x+O(x^2)$.

\section{Scope, equality and reproducibility}\label{sec:scope}
The proof determines a universal exponent and supplies a family attaining it. It does not classify every metric attaining that exponent. Nor does it give a metric-independent positive spectral gap below the endpoint: rescaling the original distances alone changes Gram eigenvalues without changing negative type.

The analytic argument uses no numerical minimisation, bounded-ratio reduction or finite covering of metric space. The packing step optimises over a bounded polytope for each fixed length vector. The inverse step selects a largest magnitude from finitely many entries of one fixed inverse matrix. Endpoint passage is pointwise in the metric and coefficient vector.

The publication supplement contains the source derivation, the supplied audit and review, exact rational checks, separate symbolic reconstructions, numerical corroboration, a claim-to-check map and cryptographic manifests. The checks verify identities and explicit examples; they do not constitute a proof-assistant verification of the analytic induction. The source and publication metadata distinguish the external small-case results from the self-contained proof of \cref{thm:main}. No priority claim is made for the classical Gram, Schur-complement, inversion or circumcentre identities.

The current numerical summary uses \texttt{null} for empty samples and reports accepted counts for every statistic. It separates centred-Gram eigenvalues, eigenvalues restricted to $\one^\perp$, and the split-vector quadratic form. Optimiser convergence and feasibility are recorded separately. Historical files are retained unchanged, including superseded sentinel values and mixed statistic names; those fields are not current estimates.

\paragraph{Research provenance.}
The argument was developed in an AI-assisted research workflow with iterative mathematical review and computational checks. This anonymous manuscript is prepared for publication by Evidence Press as an unrefereed research paper. The supplied review is preserved with a point-by-point revision response; its provenance does not establish external specialist endorsement. The release record distinguishes the completed manuscript from website publication, DOI registration and rights metadata, none of which is inferred from the proof checks.

\appendix
\section{Two inverse-pairing statements that fail}\label{app:false}
The proof avoids two superficially natural strengthenings. These examples also explain why the truncations in \eqref{eq:packing-constants} and \eqref{eq:kn} are retained.

\subsection{Shortening a star edge can increase the inverse pairing}
Take a root, four leaves of raw length one and two leaves of raw length $99$. The star distance between the long leaves is $198$. Replace this distance alone by $s\in[197,198]$ and leave all others unchanged. Every such array is a Ptolemaic metric. Triangles involving the altered edge require only $s\le198$ and $s\le200$. A quartet containing both long leaves and two short leaves has opposite-edge products $2s,10000,10000$; one containing both long leaves, the root and a short leaf has products $s,9900,9900$. All other inequalities are those of the original star.

At $p=1/2$, put
\[
 r=\sqrt{99},\quad c=1-1/\sqrt2,\quad b=(r-9)/2,
 \quad A_0=1+3c,\quad D(s)=2r-\sqrt{s}/2,
 \quad\Delta=A_0D(s)-8b^2.
\]
The constant-on-parts block equations give, at either long leaf,
\[
 u_B=\frac{A_0r-4b}{\Delta}>0,
 \qquad v_B=\frac{A_0-4b}{\Delta}<0.
\]
Here $0<b<1/2$, $D(s)>r$ and $A_0>1$, so $\Delta>r-2>0$. The contrast eigenvalues are also positive. Moreover $A_0r-4b>18-r>0$, while
\[
 A_0-4b=22-3/\sqrt2-2\sqrt{99}<0.
\]
For the last sign, squaring the positive sides reduces it to $\sqrt2>185/132$. Differentiating the inverse pairing with respect to the long-block row coefficient gives
\[
 \frac{\dd\mu}{\dd D}=-2u_Bv_B>0.
\]
Shortening $s$ increases $D(s)$ and therefore increases $\mu$. Numerically,
\[
 \mu_{s=198}\approx2.11231199339,
 \qquad \mu_{s=197}\approx2.11233712655.
\]
Thus entrywise comparison with an exact star does not imply the corresponding inequality between inverse pairings.

\subsection{The untruncated bound is false even with two leaves}
Take a star with root lengths $1,15$ at $p=1/4$. Write $r=15^{1/4}$ and $b=(r-1)/2$. Its Gram matrix and inverse pairing are
\[
 C=\begin{pmatrix}1&b\\b&r\end{pmatrix},\qquad
 \mu=\frac{2r-b(1+r)}{r-b^2}.
\]
Since $r<246/125$, exact arithmetic gives
\[
 25[2r-b(1+r)]-36[r-b^2]
 =-\tfrac12(7r^2+8r-43)>0,
\]
so $\mu>36/25$. On the other hand $T=2^{1/4}<119/100$, and therefore
\[
 R_2(T)=\frac4{4-T}<\frac{400}{281}<\frac{36}{25}.
\]
In decimals, $\mu\approx1.44167498$ whereas $R_2(T)\approx1.42308600$. The untruncated inverse-pairing bound fails. Neither example violates \cref{thm:packing}.

\section{Where the strictness hypotheses enter}\label{app:hypotheses}
The following table records the two inductions separately. It is a reading aid; each implication is proved in the cited location.

\begin{center}
\small
\begin{tabular}{@{}p{0.24\textwidth}p{0.35\textwidth}p{0.34\textwidth}@{}}
\toprule
Object or operation & Required positive-definiteness input & Source of the input \\
\midrule
Augmented inverse for $n+1$ metric points & Strict $p$-negative type of that metric & Explicit hypothesis $\stp(n+1)$ in \cref{thm:inverse-bound} \\[4pt]
Complement of two augmented indices & A regauged metric on $n-1$ points, or a direct $2\times2$ calculation & \Cref{lem:inversion,lem:pair-inertia} and $\stp(n+1)$ \\[4pt]
Delete one augmented index & A regauged metric on $n$ points & $\stp(n)$ inherited from $\stp(n+1)$ \\[4pt]
Renormalise a selected inverse pair & A possibly inverted metric on $n+1$ points & Universal quantifier in $\stp(n+1)$ \\[4pt]
One- and two-pole Schur blocks in an $m$-point metric & Metrics on $a+1$ and $a+2$ points & Strong induction: $a+2\le m-1$ \\[4pt]
Endpoint $p=p_m$ & None for an inverse at the endpoint & Pointwise continuity after strict subcritical negative type is proved \\
\bottomrule
\end{tabular}
\end{center}

For the conditional inverse induction the parameter check is
$T\le n/(n-2)$, which implies its lower-dimensional counterpart. For the cardinality induction, it follows from
\[
 T<T_m\le F(a,r)\le\frac{a+1}{a-1},\qquad n=a+1.
\]
The two uses of induction therefore do not require an embedding of the full $m$-point metric before \eqref{eq:completion} proves it.

\clearpage
\begin{thebibliography}{9}
\providecommand{\natexlab}[1]{#1}
\small
\bibitem[Baker et~al.(2026)Baker, Huh, Kummer, and Lorscheid]{bhkl}
Baker, M., Huh, J., Kummer, M., \& Lorscheid, O. (2026).
\emph{Lorentzian polynomials and matroids over triangular hyperfields 2: Analytic aspects} (Version 2) [Preprint]. arXiv.
\url{https://arxiv.org/abs/2607.15375v2}

\bibitem[Ercan(2026)]{ercan}
Ercan, E. (2026). \emph{Ptolemaic negative type and the values $q(6)$ and $q(7)$} (Version 1) [Preprint]. arXiv.
\url{https://arxiv.org/abs/2608.20606v1}

\bibitem[Fiedler(2011)]{fiedler}
Fiedler, M. (2011). \emph{Matrices and graphs in geometry}. Cambridge University Press.
\url{https://doi.org/10.1017/CBO9780511973611}

\bibitem[Haynsworth(1968)]{haynsworth}
Haynsworth, E. V. (1968). Determination of the inertia of a partitioned Hermitian matrix. \emph{Linear Algebra and its Applications, 1}(1), 73--81.
\url{https://doi.org/10.1016/0024-3795(68)90050-5}

\bibitem[Korea Superintelligence Labs(2026)]{ksl}
Korea Superintelligence Labs (2026). \emph{Sharp constants for Ptolemaic negative type: the p-anchor ladder, an impossibility window above $q=1$, and Bernstein cube inequalities on both sides of $q=1$} (Version 3, 7 September) [Preprint]. Ideosphere Machina Mathematica.
\url{https://ideosphere.ai/papers/ptolemaic-negtype-v3.pdf}

\bibitem[Li and Weston(2010)]{liweston}
Li, H., \& Weston, A. (2010). Strict p-negative type of a metric space. \emph{Positivity, 14}, 529--545.
\url{https://doi.org/10.1007/s11117-009-0035-2}

\bibitem[S\'{a}nchez(2012)]{sanchez}
S\'{a}nchez, S. (2012). On the supremal p-negative type of finite metric spaces. \emph{Journal of Mathematical Analysis and Applications, 389}(1), 98--107.
\url{https://doi.org/10.1016/j.jmaa.2011.11.041}

\bibitem[Schoenberg(1938)]{schoenberg}
Schoenberg, I. J. (1938). Metric spaces and positive definite functions. \emph{Transactions of the American Mathematical Society, 44}(3), 522--536.
\url{https://doi.org/10.1090/S0002-9947-1938-1501980-0}
\end{thebibliography}
\end{document}
